\documentclass[14pt]{extarticle}
% =============================================================================
% Shared Preamble for Worksheets, Homeworks, Exams & Lesson Notes
% (scripts/tex_preamble.tex)
% =============================================================================

\newif\ifdetailed
\ifdefined\setdetailed
  \setdetailed
\fi

\newif\ifwithsols
\ifdefined\setwithsols
  \setwithsols
\fi

\newif\ifwhiteblack
\ifdefined\setwhiteblack
  \setwhiteblack
\fi

% Document-class specific packages (loaded before common_preamble.tex / polyglossia)
\usepackage{fancyhdr}
\usepackage{enumitem}
\usepackage{titlesec}
\usepackage[bottom]{footmisc}
\usepackage{titling}
\usepackage{longtable} % page-breakable tables
\usepackage{float}
\usepackage[onehalfspacing]{setspace}

\ifwhiteblack
  \usepackage[a4paper, landscape, top=1.25cm, bottom=1.25cm, left=1.25cm, right=1.25cm, headheight=1.5cm, headsep=0.5cm]{geometry}
\else
  \usepackage[a4paper, top=2.4cm, bottom=2.4cm, left=2.2cm, right=2.2cm, headheight=1.4cm, headsep=0.4cm, footskip=1.0cm]{geometry}
\fi

% Shared common packages, fonts, polyglossia, callout boxes, code helpers, date & header macros
\IfFileExists{scripts/common_preamble.tex}{%
  % =============================================================================
% Shared Common Preamble (scripts/common_preamble.tex)
% Included by:
%   - scripts/tex_preamble.tex    (worksheets, homeworks, exams, lesson notes)
%   - scripts/book_preamble.tex   (two-sided Arabic book)
%   - scripts/beamer_preamble.tex (Beamer presentation slides)
% =============================================================================

\usepackage{fontspec}
\usepackage{amsmath}
\usepackage{amssymb}
\usepackage{xcolor}
\usepackage{graphicx}
\usepackage{listings}
\usepackage[most]{tcolorbox}
\usepackage{pifont}
\usepackage{minted}
\usepackage{etoolbox}
\usepackage{array}
\usepackage{booktabs}
\usepackage{tikz}
\usepackage{datetime2}

\usetikzlibrary{automata, positioning, arrows.meta, arrows, shapes.geometric, calc}

% Beamer already loads hyperref internally; load it here for article/book classes
% right before polyglossia (bidi).
\makeatletter
\@ifclassloaded{beamer}{}{%
  \usepackage{hyperref}%
  \hypersetup{hidelinks}%
  \providecommand{\LTcaptype}{table}%
}
\makeatother

% Load polyglossia (and bidi) AFTER graphics, tables, boxes, code, and hyperref
\usepackage{polyglossia}
\setmainlanguage[numerals=western]{arabic}
\setotherlanguage{english}
\newfontfamily\arabicfont[Script=Arabic]{Calibri}
\newfontfamily\arabicfontsf[Script=Arabic]{Calibri}
\newfontfamily\arabicfonttt[Script=Arabic]{Courier New}
\newfontfamily\emoji[BoldFont=*,ItalicFont=*]{Segoe UI Emoji}

% Allow slight inter-word stretch to avoid Overfull \hbox with inline \textenglish
\emergencystretch=1.5em

% -----------------------------------------------------------------------------
% Default Minted & Callout Box Definitions
% -----------------------------------------------------------------------------
\AtBeginEnvironment{minted}{\setstretch{0.95}}
\setminted{fontsize=\small, breaklines=true, vspace=2pt}

\tcbset{
  arc=2.5pt, outer arc=2.5pt,
  boxsep=1pt, left=6pt, right=6pt, top=3pt, bottom=3pt,
  toptitle=1.5pt, bottomtitle=1.5pt,
  before skip=6pt, after skip=6pt,
  fonttitle=\bfseries,
  fontupper=\setstretch{1.1}\setlength{\parskip}{2pt},
  colback=white, colframe=black, boxrule=0.6pt,
  boxcolor/.style 2 args={colback=#1, colframe=#2}
}

\newtcolorbox{boxDef}[1][]{boxcolor={blue!4!white}{blue!75!black},
  title={{\emoji📘} تعريف\ifx\\#1\\\else ~#1\fi :}}
\newtcolorbox{boxExercise}[1][]{boxcolor={cyan!4!white}{cyan!70!black},
  title={{\emoji🧩} تمرين\ifx\\#1\\\else ~#1\fi :}}
\newtcolorbox{boxExample}[1][]{boxcolor={yellow!6!white}{orange!85!black},
  title={{\emoji📝} مثال\ifx\\#1\\\else ~#1\fi :}}
\newtcolorbox{boxNote}[1][]{boxcolor={gray!7!white}{black!75},
  title={{\emoji✨} ملاحظة\ifx\\#1\\\else ~#1\fi :}}
\newtcolorbox{boxAttention}[1][]{boxcolor={magenta!5!white}{magenta!75!black},
  title={{\emoji🔔} تنبيه\ifx\\#1\\\else ~#1\fi :}}
\newtcolorbox{boxWarning}[1][]{boxcolor={red!4!white}{red!75!black},
  title={{\emoji⚡} ملاحظة هامة\ifx\\#1\\\else ~#1\fi :}}
\newtcolorbox{boxSolution}[1][]{boxcolor={green!4!white}{green!60!black},
  title={{\emoji✅} حل\ifx\\#1\\\else ~#1\fi :}}
\newtcolorbox{boxSymbol}[1][]{boxcolor={purple!4!white}{purple!70!black},
  title={{\emoji🔣} رمز\ifx\\#1\\\else ~#1\fi :}}
\newtcolorbox{boxHint}[1][]{boxcolor={teal!4!white}{teal!65!black},
  title={{\emoji💡} تلميح\ifx\\#1\\\else ~#1\fi :}}

% Compact code boxes: tighter line spacing and small vertical gaps
\tcbset{
  simplecode/.style={
    arc=2pt, outer arc=2pt,
    left=4pt, right=4pt, top=4pt, bottom=4pt,
    fontupper=\setstretch{0.95},
    before skip=6pt, after skip=6pt,
    colback=gray!5, colframe=black!50, boxrule=0.4pt
  }
}
\newenvironment{boxCode}{\begin{tcolorbox}[simplecode]}{\end{tcolorbox}}

% Code cell used inside tables (keeps code left-to-right inside RTL tables)
% Use inside a fixed-width p{...} column: \begin{pycell}\begin{minted}...\end{minted}\end{pycell}
\makeatletter
\newtcolorbox{pycell}{frame empty,colback=white,boxsep=0pt,left=2pt,right=2pt,top=2pt,bottom=2pt,width=\linewidth,before=\noindent,after=,baseline={\tcb@height-\ht\strutbox},fontupper=\setstretch{0.95}\setlength{\topsep}{0pt}\setlength{\partopsep}{0pt}\setlength{\parskip}{0pt}}
\makeatother

\newcolumntype{C}[1]{>{\centering\arraybackslash}p{#1}}

% Dotted date style (DD.MM.YYYY)
\makeatletter
\DTMnewdatestyle{dotted}{%
  \renewcommand{\DTMdisplaydate}[4]{\two@digits{##3}.\two@digits{##2}.##1}%
  \renewcommand{\DTMDisplaydate}{\DTMdisplaydate}%
}
\makeatother
\DTMsetdatestyle{dotted}

% -----------------------------------------------------------------------------
% Customizable Header / Footer Metadata
% -----------------------------------------------------------------------------
\newcommand{\headerleft}{مدرسة التسامح الشاملة}
\newcommand{\headerright}{الأستاذ محمود اغبارية}
\newcommand{\footercenter}{\thepage}

\newcommand{\setheaderleft}[1]{\renewcommand{\headerleft}{#1}}
\newcommand{\setheaderright}[1]{\renewcommand{\headerright}{#1}}
\newcommand{\setfootercenter}[1]{\renewcommand{\footercenter}{#1}}
%
}{%
  \IfFileExists{../../../scripts/common_preamble.tex}{%
    % =============================================================================
% Shared Common Preamble (scripts/common_preamble.tex)
% Included by:
%   - scripts/tex_preamble.tex    (worksheets, homeworks, exams, lesson notes)
%   - scripts/book_preamble.tex   (two-sided Arabic book)
%   - scripts/beamer_preamble.tex (Beamer presentation slides)
% =============================================================================

\usepackage{fontspec}
\usepackage{amsmath}
\usepackage{amssymb}
\usepackage{xcolor}
\usepackage{graphicx}
\usepackage{listings}
\usepackage[most]{tcolorbox}
\usepackage{pifont}
\usepackage{minted}
\usepackage{etoolbox}
\usepackage{array}
\usepackage{booktabs}
\usepackage{tikz}
\usepackage{datetime2}

\usetikzlibrary{automata, positioning, arrows.meta, arrows, shapes.geometric, calc}

% Beamer already loads hyperref internally; load it here for article/book classes
% right before polyglossia (bidi).
\makeatletter
\@ifclassloaded{beamer}{}{%
  \usepackage{hyperref}%
  \hypersetup{hidelinks}%
  \providecommand{\LTcaptype}{table}%
}
\makeatother

% Load polyglossia (and bidi) AFTER graphics, tables, boxes, code, and hyperref
\usepackage{polyglossia}
\setmainlanguage[numerals=western]{arabic}
\setotherlanguage{english}
\newfontfamily\arabicfont[Script=Arabic]{Calibri}
\newfontfamily\arabicfontsf[Script=Arabic]{Calibri}
\newfontfamily\arabicfonttt[Script=Arabic]{Courier New}
\newfontfamily\emoji[BoldFont=*,ItalicFont=*]{Segoe UI Emoji}

% Allow slight inter-word stretch to avoid Overfull \hbox with inline \textenglish
\emergencystretch=1.5em

% -----------------------------------------------------------------------------
% Default Minted & Callout Box Definitions
% -----------------------------------------------------------------------------
\AtBeginEnvironment{minted}{\setstretch{0.95}}
\setminted{fontsize=\small, breaklines=true, vspace=2pt}

\tcbset{
  arc=2.5pt, outer arc=2.5pt,
  boxsep=1pt, left=6pt, right=6pt, top=3pt, bottom=3pt,
  toptitle=1.5pt, bottomtitle=1.5pt,
  before skip=6pt, after skip=6pt,
  fonttitle=\bfseries,
  fontupper=\setstretch{1.1}\setlength{\parskip}{2pt},
  colback=white, colframe=black, boxrule=0.6pt,
  boxcolor/.style 2 args={colback=#1, colframe=#2}
}

\newtcolorbox{boxDef}[1][]{boxcolor={blue!4!white}{blue!75!black},
  title={{\emoji📘} تعريف\ifx\\#1\\\else ~#1\fi :}}
\newtcolorbox{boxExercise}[1][]{boxcolor={cyan!4!white}{cyan!70!black},
  title={{\emoji🧩} تمرين\ifx\\#1\\\else ~#1\fi :}}
\newtcolorbox{boxExample}[1][]{boxcolor={yellow!6!white}{orange!85!black},
  title={{\emoji📝} مثال\ifx\\#1\\\else ~#1\fi :}}
\newtcolorbox{boxNote}[1][]{boxcolor={gray!7!white}{black!75},
  title={{\emoji✨} ملاحظة\ifx\\#1\\\else ~#1\fi :}}
\newtcolorbox{boxAttention}[1][]{boxcolor={magenta!5!white}{magenta!75!black},
  title={{\emoji🔔} تنبيه\ifx\\#1\\\else ~#1\fi :}}
\newtcolorbox{boxWarning}[1][]{boxcolor={red!4!white}{red!75!black},
  title={{\emoji⚡} ملاحظة هامة\ifx\\#1\\\else ~#1\fi :}}
\newtcolorbox{boxSolution}[1][]{boxcolor={green!4!white}{green!60!black},
  title={{\emoji✅} حل\ifx\\#1\\\else ~#1\fi :}}
\newtcolorbox{boxSymbol}[1][]{boxcolor={purple!4!white}{purple!70!black},
  title={{\emoji🔣} رمز\ifx\\#1\\\else ~#1\fi :}}
\newtcolorbox{boxHint}[1][]{boxcolor={teal!4!white}{teal!65!black},
  title={{\emoji💡} تلميح\ifx\\#1\\\else ~#1\fi :}}

% Compact code boxes: tighter line spacing and small vertical gaps
\tcbset{
  simplecode/.style={
    arc=2pt, outer arc=2pt,
    left=4pt, right=4pt, top=4pt, bottom=4pt,
    fontupper=\setstretch{0.95},
    before skip=6pt, after skip=6pt,
    colback=gray!5, colframe=black!50, boxrule=0.4pt
  }
}
\newenvironment{boxCode}{\begin{tcolorbox}[simplecode]}{\end{tcolorbox}}

% Code cell used inside tables (keeps code left-to-right inside RTL tables)
% Use inside a fixed-width p{...} column: \begin{pycell}\begin{minted}...\end{minted}\end{pycell}
\makeatletter
\newtcolorbox{pycell}{frame empty,colback=white,boxsep=0pt,left=2pt,right=2pt,top=2pt,bottom=2pt,width=\linewidth,before=\noindent,after=,baseline={\tcb@height-\ht\strutbox},fontupper=\setstretch{0.95}\setlength{\topsep}{0pt}\setlength{\partopsep}{0pt}\setlength{\parskip}{0pt}}
\makeatother

\newcolumntype{C}[1]{>{\centering\arraybackslash}p{#1}}

% Dotted date style (DD.MM.YYYY)
\makeatletter
\DTMnewdatestyle{dotted}{%
  \renewcommand{\DTMdisplaydate}[4]{\two@digits{##3}.\two@digits{##2}.##1}%
  \renewcommand{\DTMDisplaydate}{\DTMdisplaydate}%
}
\makeatother
\DTMsetdatestyle{dotted}

% -----------------------------------------------------------------------------
% Customizable Header / Footer Metadata
% -----------------------------------------------------------------------------
\newcommand{\headerleft}{مدرسة التسامح الشاملة}
\newcommand{\headerright}{الأستاذ محمود اغبارية}
\newcommand{\footercenter}{\thepage}

\newcommand{\setheaderleft}[1]{\renewcommand{\headerleft}{#1}}
\newcommand{\setheaderright}[1]{\renewcommand{\headerright}{#1}}
\newcommand{\setfootercenter}[1]{\renewcommand{\footercenter}{#1}}
%
  }{%
    \input{../../scripts/common_preamble.tex}%
  }%
}

\pagestyle{fancy}

% Prevent widow/orphan lines
\widowpenalty=10000
\clubpenalty=10000

% Compact section headings and default list spacing so pages do not waste vertical space
\titlespacing*{\section}{0pt}{1.4ex plus 0.4ex minus 0.2ex}{0.7ex plus 0.2ex}
\titlespacing*{\subsection}{0pt}{1.1ex plus 0.3ex minus 0.2ex}{0.5ex plus 0.1ex}
\titlespacing*{\subsubsection}{0pt}{0.9ex plus 0.2ex minus 0.1ex}{0.4ex plus 0.1ex}
\setlist{topsep=4pt, partopsep=0pt, parsep=2pt, itemsep=4pt}

\gappto\captionsarabic{\renewcommand{\contentsname}{جدول المحتويات}}

% Legacy listings settings
\lstset{
  language=[Sharp]C,
  numbers=left,
  stepnumber=1,
  numbersep=8pt,
  frame=single,
  basicstyle=\ttfamily\small
}
\ifwhiteblack
  \lstset{
    keywordstyle=\bfseries,
    stringstyle=\ttfamily,
    commentstyle=\itshape
  }
\else
  \lstset{
    keywordstyle=\color{blue},
    stringstyle=\color{red},
    commentstyle=\color{green!50!black}
  }
\fi

% White-and-black landscape presentation mode overrides
\ifwhiteblack
  \AtBeginDocument{
    \LARGE
    \usemintedstyle{bw}
  }
  \setminted{vspace=2pt}
  \tcbset{
    colback=white, colframe=black, boxrule=1pt,
    boxcolor/.style 2 args={colback=white, colframe=black},
    simplecode/.append style={colback=white, colframe=black, boxrule=1pt}
  }
\fi

% Default to empty author to avoid "LaTeX Warning: No \author given" when omitted
\author{}

% Compact title block
\makeatletter
\renewcommand{\@maketitle}{%
  \begin{center}%
    {\ifwhiteblack\Huge\else\LARGE\fi \bfseries \@title \par}%
    \ifdefempty{\@author}{}{%
      \vskip 0.15em%
      {\large \@author \par}%
    }%
    {\small \hfill \textenglish{\DTMtoday} \par}%
  \end{center}%
  \vspace{-0.3cm}%
}
\makeatother

\apptocmd{\maketitle}{\thispagestyle{fancy}}{}{}

% Apply page headers and footer (using \headerleft, \headerright, \footercenter from common_preamble.tex)
\newcommand{\applyfancyheaders}{%
  \fancyhf{}%
  \fancyhead[L]{\textarabic{\headerleft}}%
  \fancyhead[R]{\textarabic{\headerright}}%
  \fancyfoot[C]{\textarabic{\footercenter}}%
}

\fancypagestyle{plain}{\applyfancyheaders}
\applyfancyheaders

\setcounter{tocdepth}{3} % only section, subsection, and subsubsection in TOC

\newcommand{\importpdfpage}[2]{%
    \newgeometry{left=0cm,right=0cm,top=0cm,bottom=0cm}
    \thispagestyle{empty}
    \includegraphics[page=#2,width=0.95\paperwidth,height=0.95\paperheight,keepaspectratio]{#1}
    \restoregeometry
}

\newcommand{\insertFullImg}[1]{%
    \noindent
    \makebox[\textwidth][c]{
        \includegraphics[width=0.9\paperwidth,keepaspectratio]{#1}%
    }%
}%


\title{ملخص: العمليات الخارجية في لغة \textenglish{C\#}}

\begin{document}
\maketitle
\thispagestyle{fancy}

\section*{مقدمة ومصطلحات أساسية}
العمليات الخارجية تقسّم البرنامج إلى مقاطع صغيرة ومستقلة تؤدي مهام محددة، مما يسهّل قراءة الكود وصيانته ويمنع تكراره.

\begin{boxAttention}[- مصطلحات أساسية]
\begin{itemize}[nosep]
    \item \textbf{تستقبل:} تقرأ قيمة من المستخدم عبر لوحة المفاتيح باستخدام \textenglish{Console.ReadLine()}.
    \item \textbf{تتلقى:} تأخذ معطيات (بارامترات) تُمرَّر إليها بين القوسين في تعريف العملية وعند استدعائها.
    \item \textbf{تطبع:} تعرض مخرجات على الشاشة باستخدام \textenglish{Console.WriteLine(...)}.
    \item \textbf{تعيد:} تُرجع قيمة إلى المكان الذي استُدعيت منه باستخدام الأمر \textenglish{return}.
\end{itemize}
\end{boxAttention}

\subsection*{ما هي الفئة (\textenglish{class}) وأين تعيش العمليات؟}
\textbf{الفئة (\textenglish{class})} في لغة \textenglish{C\#} هي الإطار (أو الحاوية) الذي يضمّ الكود البرمجي. جميع العمليات الخارجية — بما فيها العملية الرئيسية \textenglish{Main} — \textbf{تعيش داخل فئة} (مثل \textenglish{class Program})، وتُكتب الواحدة تلو الأخرى داخل قوسي الفئة (وليس عملية داخل قوسي عملية أخرى).

\subsection*{المبنى العام لتعريف العملية الخارجية داخل الفئة}
\begin{boxCode}
\begin{english}
\begin{minted}{csharp}
class Program
{
    public static returnType MethodName(parameters)
    {
        // Method body (statements)
    }

    public static void Main(string[] args)
    {
        // Calling the method here
    }
}
\end{minted}
\end{english}
\end{boxCode}

\begin{itemize}[nosep]
    \item \textenglish{public static}: كلمتان ثابتتان في بداية تعريف كل عملية خارجية في هذه المرحلة.
    \item \textenglish{returnType}: نمط القيمة التي تعيدها العملية (\textenglish{int, double, string, char, bool})، أو \textenglish{void} إذا كانت لا تعيد قيمة.
    \item \textenglish{MethodName}: اسم العملية (يبدأ بحرف كبير \textenglish{PascalCase}).
    \item \textenglish{parameters}: المعطيات التي تتلقاها العملية مع أنماطها، أو قوسان فارغان \textenglish{()} إذا كانت لا تتلقى معطيات.
\end{itemize}

% ==============================================================
% القسم الأول: عمليات لا تتلقى ولا تعيد
% ==============================================================
\filbreak
\section{عمليات لا تتلقى ولا تعيد}
نمط الإرجاع فيها \textenglish{void} وأقواسها فارغة \textenglish{()}. عند استدعائها في \textenglish{Main} نكتب اسمها فقط دون حفظها في متغير.

\begin{boxExample}[عملية تستقبل عددًا $n$ وتطبع $n$ نجوم باستخدام حلقة]
\begin{english}
\begin{minted}{csharp}
public static void ReadAndPrintStars()
{
    int n = int.Parse(Console.ReadLine());
    for (int i = 0; i < n; i++)
    {
        Console.WriteLine("*");
    }
}

public static void Main()
{
    ReadAndPrintStars();
}
\end{minted}
\end{english}
\end{boxExample}

% ==============================================================
% القسم الثاني: عمليات تتلقى ولا تعيد
% ==============================================================
\filbreak
\section{عمليات تتلقى ولا تعيد}
نمط الإرجاع فيها \textenglish{void}، وتأخذ معطيات بين القوسين. عند استدعائها نمرّر لها قيمًا تطابق المعطيات في النمط والعدد والترتيب (ويمكن أن تستقبل مدخلات إضافية داخلها).

\begin{boxExample}[عملية تتلقى أجر الساعة وتستقبل عدد الساعات ثم تطبع الراتب]
\begin{english}
\begin{minted}{csharp}
public static void PrintSalary(double hourlyRate)
{
    int hours = int.Parse(Console.ReadLine());
    double salary = hourlyRate * hours;
    Console.WriteLine("Salary = " + salary);
}

public static void Main()
{
    PrintSalary(45.0);
}
\end{minted}
\end{english}
\end{boxExample}

% ==============================================================
% القسم الثالث: عمليات لا تتلقى وتعيد أعدادًا صحيحة أو عشرية
% ==============================================================
\filbreak
\section{عمليات لا تتلقى وتعيد أعدادًا (\textenglish{int} أو \textenglish{double})}
أقواسها فارغة \textenglish{()}، وتعيد نتيجة عددية بالأمر \textenglish{return}. عند استدعائها في \textenglish{Main} نحفظ القيمة المعادة في متغير من نفس النمط أو نطبعها مباشرة.

\begin{boxExample}[عملية تستقبل $5$ أسعار بحلقة \textenglish{for} وتعيد مجموعها]
\begin{english}
\begin{minted}{csharp}
public static double ReadFivePricesSum()
{
    double sum = 0;
    for (int i = 0; i < 5; i++)
    {
        double price = double.Parse(Console.ReadLine());
        sum += price;
    }
    return sum;
}

public static void Main()
{
    double total = ReadFivePricesSum();
    Console.WriteLine("Total = " + total);
}
\end{minted}
\end{english}
\end{boxExample}

% ==============================================================
% القسم الرابع: عمليات تتلقى وتعيد أعدادًا صحيحة أو عشرية
% ==============================================================
\filbreak
\section{عمليات تتلقى وتعيد أعدادًا (\textenglish{int} أو \textenglish{double})}
تتلقى معطيات بين القوسين، وتُجري حسابات عليها (بما في ذلك الشروط أو الحلقات)، ثم تعيد الناتج بالأمر \textenglish{return}.

\begin{boxExample}[عملية تتلقى عددًا $n$ وتعيد مجموع مضاعفات الـ $3$ من $1$ حتى $n$]
\begin{english}
\begin{minted}{csharp}
public static int SumMultiplesOf3(int n)
{
    int sum = 0;
    for (int i = 1; i <= n; i++)
    {
        if (i % 3 == 0)
        {
            sum += i;
        }
    }
    return sum;
}

public static void Main()
{
    int res = SumMultiplesOf3(15);
    Console.WriteLine(res);
}
\end{minted}
\end{english}
\end{boxExample}

% ==============================================================
% القسم الخامس: عمليات تعيد string أو char
% ==============================================================
\filbreak
\section{عمليات تعيد \textenglish{string} أو \textenglish{char}}
تعيد نصًّا (\textenglish{string}) بين علامتي اقتباس مزدوجتين \textenglish{"..."} أو حرفًا واحدًا (\textenglish{char}) بين علامتي اقتباس مفردتين \textenglish{'...'}.

\begin{boxExample}[عملية تعيد نصًّا (\textenglish{string}) وأخرى تعيد رمز الإشارة (\textenglish{char})]
\begin{english}
\begin{minted}{csharp}
public static string Greet(string name)
{
    return "Hello " + name;
}

public static char GetSignSymbol(int num)
{
    if (num >= 0)
    {
        return '+';
    }
    else
    {
        return '-';
    }
}

public static void Main()
{
    Console.WriteLine(Greet("Ali"));
    char sym = GetSignSymbol(-12);
    Console.WriteLine(sym);
}
\end{minted}
\end{english}
\end{boxExample}

% ==============================================================
% القسم السادس: عمليات تعيد bool وكيفية استدعائها
% ==============================================================
\filbreak
\section{عمليات تعيد \textenglish{bool} واستدعاؤها داخل \textenglish{if}}
تعيد \textenglish{true} أو \textenglish{false}، ويمكن استدعاؤها مباشرة داخل قوسي جملة \textenglish{if (...)}.

\begin{boxExample}[تعريف عملية تعيد \textenglish{bool} بطريقتين واستدعاؤها في \textenglish{if}]
\begin{english}
\begin{minted}{csharp}
// Way 1: using if - else
public static bool IsPositive(int num)
{
    if (num > 0)
    {
        return true;
    }
    else
    {
        return false;
    }
}

// Way 2: returning the condition directly
public static bool IsAdult(int age)
{
    return (age >= 18);
}

public static void Main()
{
    int n = int.Parse(Console.ReadLine());
    if (IsPositive(n))
    {
        Console.WriteLine("Positive");
    }
    else
    {
        Console.WriteLine("Zero or negative");
    }
}
\end{minted}
\end{english}
\end{boxExample}

% ==============================================================
% القسم السابع: عمليات تستدعي عمليات أخرى
% ==============================================================
\filbreak
\section{عمليات تستدعي عمليات أخرى}
يمكن لأي عملية خارجية أن تستدعي عملية أخرى لمساعدتها في حساب النتيجة.

\begin{boxExample}[عملية تحسب تكعيب عدد باستدعاء عملية الضرب]
\begin{english}
\begin{minted}{csharp}
public static int Multiply(int a, int b)
{
    return a * b;
}

public static int Cube(int x)
{
    int temp = Multiply(x, x);
    return Multiply(temp, x);
}

public static void Main()
{
    Console.WriteLine(Cube(4));
}
\end{minted}
\end{english}
\end{boxExample}

% ==============================================================
% القسم الثامن: ملاحظات واستخدام عمليات جاهزة
% ==============================================================
\filbreak
\section{استدعاء عمليات جاهزة وملاحظات عامة}
\begin{itemize}[nosep]
    \item \textbf{استدعاء عملية معطى تعريفها فقط:} في بعض الأسئلة يُعطى تعريف العملية ووظيفتها فقط دون كودها الداخلي؛ عندها نستدعيها مباشرة حسب تعريفها دون إعادة كتابتها.
    \item \textbf{استدعاء عملية من فئة أخرى:} نكتب اسم الفئة ثم نقطة ثم اسم العملية، مثل: \\
    \textenglish{Console.WriteLine(...)}، \textenglish{int.Parse(...)}، وعمليات الفئة \textenglish{Math} (\textenglish{Math.Max}، \textenglish{Math.Abs}، \textenglish{Math.Sqrt}).
    \item \textbf{دمج الاستدعاء في تعبير أو حلقة:} يمكن استدعاء العملية التي تعيد قيمة داخل تعبير حسابي أو شرط أو داخل حلقة \textenglish{for} / \textenglish{while}، أو تمرير نتيجتها كمعطى لعملية أخرى مثل \textenglish{Math.Max(a, Math.Max(b, c))}.
\end{itemize}

\vspace{1cm}
\begin{flushleft}
أرجو لكم وقتًا ممتعًا.

الأستاذ محمود اغبارية.
\end{flushleft}

\end{document}
